%% kapa: exam=1 round=30 year=2019 period=1 form=A subject=economics subject_ko=경제학원론 seq=24 number=64
%% source: Q-Net 감정평가사 기출문제 (cst00309) artlSeq=5205205 file=extracted/1008/5205205/2019년 제30회 감정평가사 1차시험 문제지/2019년 제30회 감정평가사 1차시험 문제지/2019년도 제30회 감정평가사 1차 1교시 A형.hwp
%% conversion: kordoc 4.24.51 -> tools/gen1.py (auto); review=proofread-sonnet; flags=none
\begin{question}{64}
개방경제인 甲국의 국민소득 결정모형이 다음과 같을 때, 甲국의 국내총소득, 국민총소득, 처분가능소득은? (단, 제시된 항목 외 다른 것은 고려하지 않는다.)
\begin{framedtext}
○ 국내총생산: 1,000 \par
○ 대외 순수취 요소소득: 20 \par
○ 교역조건 변화에 따른 실질무역 손익: 50 \par
○ 감가상각: 10 \par
○ 사내유보이윤: 10 \par
○ 각종세금: 3 \par
○ 이전지출: 3
\end{framedtext}
\begin{choices}
  \item 1,000, 980, 960
  \item 1,000, 1,020, 1,000
  \item 1,050, 1,050, 1,050
  \item 1,050, 1,070, 1,050
  \item 1,070, 1,050, 1,030
\end{choices}
\end{question}
\answer{4}
